%=============================================================================
\section{传感器标定与选型}
%=============================================================================

上一章的滤波器能去噪，但修不了一个在骗你的传感器。陀螺仪有 $2\textdegree{}/\mathrm{s}$ 的偏置（Bias），不管低通滤波器多厉害，一分钟照样漂 $120\textdegree{}$。这一章教你在把数据喂进控制流水线之前，先让传感器说实话。

标定这活儿不炫酷——无非是某个周六下午，搬张平桌、拿个量角器，蹲在那里采数据。但它是机器人开发中性价比最高的操作：花十分钟标定，能消灭那些多少行代码都补不回来的误差。

%-----------------------------------------------------------------------------
\subsection{传感器误差模型}
%-----------------------------------------------------------------------------

你用到的每一个传感器，误差模型大同小异。测量值 $y_m$ 和真值 $y$ 之间的关系是：

\begin{equation}
y_m = S \, y + b + n
\end{equation}

其中 $S$ 是\textbf{比例因子}（Scale Factor，理想值为 1，实际上 MEMS 传感器一般在 0.98--1.02 之间），$b$ 是\textbf{偏置}（Bias，即使真值为零也会存在的常数偏移），$n$ 是\textbf{噪声}（Noise，零均值的随机波动）。

来逐个建立直觉。\textbf{偏置}是最阴险的：如果陀螺仪在机器人纹丝不动时读出 $0.5\textdegree{}/\mathrm{s}$，这半度每秒会不停地往上累加。一分钟后你就有了 $30\textdegree{}$ 的幽灵旋转，五分钟后 $150\textdegree{}$。没有任何滤波器能从速率传感器中去掉一个直流偏移——你必须测量它，然后减掉。

\textbf{比例因子}误差更隐蔽。如果 $S = 1.03$ 而非 $1.0$，一个 $90\textdegree{}$ 的转弯会被读成 $92.7\textdegree{}$。短暂测试你不会注意到，但打完一场比赛就累积起来了。

\textbf{噪声} $n$ 是你已经会对付的东西：低通滤波、互补滤波、卡尔曼滤波。三种误差里最好处理的。

对于多轴传感器（加速度计、陀螺仪、磁力计），模型推广为：

\begin{equation}
\mathbf{y}_m = \mathbf{S} \, \mathbf{y} + \mathbf{b} + \mathbf{n}
\end{equation}

此时 $\mathbf{S}$ 变成了一个 $3 \times 3$ 矩阵。$\mathbf{S}$ 的非对角元素体现的是\textbf{交叉轴耦合}（Cross-axis Coupling）——$x$ 轴加速度计因芯片内部机械对准不完美，会捎带上一点 $y$ 轴的加速度。对于廉价 MEMS 传感器，交叉轴耦合一般是主轴的 1--3\%，不大但可以测出来。

\begin{center}
\begin{tabular}{llll}
\toprule
\textbf{误差来源} & \textbf{典型量级} & \textbf{检测方法} & \textbf{修正方法} \\
\midrule
偏置（加速度计） & $\pm 50\;\mathrm{mg}$ & 静态测试：静止读数 & 减去均值 \\
偏置（陀螺仪） & $\pm 1\textdegree{}/\mathrm{s}$ & 静态测试：静止读数 & 减去均值 \\
比例因子 & 1--3\% & 与已知参考比较 & 求解标定方程 \\
交叉轴耦合 & 1--3\% & 六面标定 & 完整 $3\times3$ 矩阵 \\
噪声（加速度计） & $\sim\!200\;\mu g/\!\sqrt{\mathrm{Hz}}$ & Allan 方差或功率谱密度 & 低通滤波 \\
噪声（陀螺仪） & $\sim\!0.01\textdegree{}/\mathrm{s}/\!\sqrt{\mathrm{Hz}}$ & Allan 方差或功率谱密度 & 低通滤波 \\
温度漂移 & 视型号而定 & 不同温度下测试 & 在线补偿 \\
\bottomrule
\end{tabular}
\end{center}

温度漂移值得一提。MEMS 传感器的偏置会随温度变化——有时变化很大。高端 IMU 内置温度传感器，用多项式做在线补偿。对于在体育馆里打比赛的机器人来说，温度基本恒定，通常可以忽略；但要注意，冬天在室内标好的机器人，夏天搬到室外可能会有明显偏差。

%-----------------------------------------------------------------------------
\subsection{IMU 标定}
%-----------------------------------------------------------------------------

\subsubsection{加速度计：六面静态标定}

原理简洁到优雅：重力给了你一个免费的、已知的、高精度参考向量。地球表面 $\|\mathbf{g}\| = 9.81\;\mathrm{m/s^2}$，精确到四位有效数字。只要让加速度计的每根轴依次正向、反向对准重力方向，就能得到六个数据点，完全约束住 $3 \times 3$ 的比例-交叉轴矩阵 $\mathbf{S}$ 和 $3 \times 1$ 的偏置向量 $\mathbf{b}$。

把 IMU 放在平面上，摆出六个朝向：$+x$ 朝上、$-x$ 朝上、$+y$ 朝上、$-y$ 朝上、$+z$ 朝上、$-z$ 朝上。每个朝向采集几百个样本，取平均。不需要精密夹具——一张平桌加一个棱角分明的盒子就够了。

在每个朝向 $i$ 下，测量模型为：

\begin{equation}
\mathbf{a}_{m,i} = \mathbf{S} \, \mathbf{g}_i + \mathbf{b}
\end{equation}

其中 $\mathbf{g}_i$ 是该朝向下的已知重力向量。比如 $+z$ 朝上时，$\mathbf{g}_i = [0, 0, 9.81]^T$；$-x$ 朝上时，$\mathbf{g}_i = [-9.81, 0, 0]^T$。

把六组方程堆成一个矩阵方程。定义 $\mathbf{A}_m \in \mathbb{R}^{6 \times 3}$ 为测量加速度矩阵（每行一个朝向），$\mathbf{G} \in \mathbb{R}^{6 \times 3}$ 为对应的已知重力向量。则：

\begin{equation}
\mathbf{A}_m = \mathbf{G} \, \mathbf{S}^T + \mathbf{1} \, \mathbf{b}^T
\end{equation}

这是一个标准线性系统。给 $\mathbf{G}$ 补一列全1来吸收偏置，然后用最小二乘求解：

\begin{equation}
\begin{bmatrix} \mathbf{S}^T \\ \mathbf{b}^T \end{bmatrix}
= \left( \tilde{\mathbf{G}}^T \tilde{\mathbf{G}} \right)^{-1} \tilde{\mathbf{G}}^T \mathbf{A}_m
\end{equation}

其中 $\tilde{\mathbf{G}} = [\mathbf{G} \;\; \mathbf{1}]$ 是增广矩阵。有了 $\mathbf{S}$ 和 $\mathbf{b}$，校正后的读数为：

\begin{equation}
\mathbf{a}_{\text{true}} = \mathbf{S}^{-1} (\mathbf{a}_m - \mathbf{b})
\end{equation}

下面是一个完成整个标定过程的 Python 脚本：

\begin{verbatim}
import numpy as np

# a_meas: (6, 3) array of mean accelerometer readings per orientation
# g_ref:  (6, 3) array of known gravity vectors per orientation
G_aug = np.hstack([g_ref, np.ones((6, 1))])          # (6, 4)
params, _, _, _ = np.linalg.lstsq(G_aug, a_meas, rcond=None)
S = params[:3, :].T                                   # (3, 3)
b = params[3, :]                                       # (3,)
S_inv = np.linalg.inv(S)

def calibrate(a_raw):
    return S_inv @ (a_raw - b)
\end{verbatim}

标定完成后，验证一下：把 IMU 转到任意朝向，检查 $\|\mathbf{a}_{\text{true}}\| \approx 9.81\;\mathrm{m/s^2}$。如果模值偏差超过 $\pm 0.05\;\mathrm{m/s^2}$，说明哪里出了问题——大概率是六个朝向中有某个没放平。

\subsubsection{陀螺仪：静态偏置与噪声}

陀螺仪的标定比加速度计简单，因为只需要一个测试条件：\textit{完全静止}。

把 IMU 放在稳固的平面上，采集 60 秒数据。每根轴的平均读数就是偏置，在代码里减掉即可：

\begin{verbatim}
# During initialization (robot stationary):
gyro_bias = np.mean(gyro_samples, axis=0)    # average over ~60s
# During operation:
gyro_corrected = gyro_raw - gyro_bias
\end{verbatim}

对大多数比赛应用来说，这就是全部标定了。每次上电都做一次，因为 MEMS 陀螺仪的偏置每次开机都会微微漂移（``开机偏置''，Turn-on Bias）。

如果想更深入地了解传感器的噪声结构，\textbf{Allan 方差}（Allan Variance）可以告诉你更多。Allan 方差的做法是：对陀螺仪输出按越来越长的时间窗口 $\tau$ 取平均，然后在双对数坐标中画出标准差 $\sigma(\tau)$ 与 $\tau$ 的关系。

Allan 方差图的关键特征如下：

在短平均时间（$\tau < 1\;\mathrm{s}$）段，曲线以 $-1/2$ 斜率下降，这里的主导噪声是\textbf{角度随机游走}（Angle Random Walk，白噪声）——传感器的基础噪声底。在 $\tau = 1\;\mathrm{s}$ 处的 $\sigma$ 值直接给出角度随机游走系数。

在中间平均时间段，曲线达到最低点。这个最低点就是\textbf{偏置不稳定性}（Bias Instability）——传感器能达到的最佳稳定度。对一颗两块钱的 MEMS 陀螺仪，一般在 $1$--$10\textdegree{}/\mathrm{hr}$ 左右；光纤陀螺仪可以到 $0.001\textdegree{}/\mathrm{hr}$。

在长平均时间（$\tau > 100\;\mathrm{s}$）段，曲线以 $+1/2$ 斜率上升，这是\textbf{速率随机游走}（Rate Random Walk）——温度变化等慢漂移过程的体现。

对于比赛机器人，别在这上面想太多。测个静态偏置、减掉、走人。Allan 方差分析在选型或调试长期漂移时有用，但你的比赛每局只有 3--5 分钟，简单减偏置完全够用。

\subsubsection{陀螺仪比例因子}

如果你还需要标定陀螺仪的比例因子（不仅仅是偏置），那就需要一个已知的旋转速率。穷人版方案：把 IMU 固定在转盘上（懒人转盘就行），精确转 $360\textdegree{}$，看积分后的陀螺仪输出是否吻合。如果陀螺仪说你转了 $354\textdegree{}$，那比例因子就是 $360/354 \approx 1.017$。对比赛机器人来说这一步很少必要——MEMS 陀螺仪的比例因子通常在标称值的 1\% 以内——但花五分钟就能多消灭一个误差源。

更精确的方法是用速率转台，以已知恒定角速度旋转。在多个转速下比较陀螺仪读数与已知速率，拟合直线，斜率就是比例因子，截距就是偏置。对大多数比赛场景来说这过于奢侈了，不过传感器厂商就是这么标定的。

%-----------------------------------------------------------------------------
\subsection{编码器：分辨率与噪声的博弈}
%-----------------------------------------------------------------------------

\subsubsection{从位置估计速度}

编码器给你的是位置（角度），但控制回路几乎总是需要速度。最简单的做法是差分：

\begin{equation}
v[n] = \frac{\theta[n] - \theta[n-1]}{T_s}
\end{equation}

其中 $T_s$ 是采样周期。高速时效果不错，但低速时因为\textbf{量化}（Quantization）就垮了。一个每转 $N$ 个计数的编码器，最小可检测位置变化为 $\Delta\theta = 2\pi / N$ 弧度，因此能测到的最小非零速度为：

\begin{equation}
v_{\min} = \frac{2\pi / N}{T_s}
\end{equation}

代数字进去看看。1024 线编码器，$1\;\mathrm{kHz}$ 采样，$v_{\min} = 2\pi / (1024 \times 0.001) \approx 6.14\;\mathrm{rad/s} \approx 58.6\;\mathrm{RPM}$。低于这个速度，速度估计值在 0 和 $v_{\min}$ 之间跳来跳去，中间什么都没有。电机明明在平稳转，却看起来像在抽搐。

解决方案按复杂度递增排列如下。

\textbf{降低速度计算的采样率。}如果每 $10\;\mathrm{ms}$ 算一次速度而不是每 $1\;\mathrm{ms}$，$v_{\min}$ 就降低 10 倍。代价是速度更新慢了 10 倍，对内环来说可能太慢了。

\textbf{M/T 方法。}不再数固定时间内的边沿数（M 方法），而是测量两个边沿之间的时间（T 方法），或者两者结合（M/T 方法）。低速时边沿到来很慢，测量边沿间隔时间可以提供精细得多的速度分辨率。这正是硬件定时器捕获/比较外设的设计用途。在 STM32 上，你把定时器通道配置为``输入捕获''模式：每次编码器边沿到来，硬件会锁存定时器计数值。两次捕获的时间差除以每个边沿对应的角度，就得到了速度，分辨率仅受限于定时器时钟（一般 $72$--$168\;\mathrm{MHz}$）。M/T 混合方法在高速时自动使用边沿计数，低速时自动使用时间测量，在整个速度范围内两全其美。

\textbf{对差分做低通滤波。}给速度估计值加一个简单的一阶低通滤波器，就能平滑掉量化跳变。会引入少量延迟，但在简洁性和效果之间往往是最好的折中：

\begin{verbatim}
v_raw = (theta - theta_prev) / Ts;
v_filtered += alpha * (v_raw - v_filtered);
theta_prev = theta;
\end{verbatim}

\textbf{基于观测器的估计。}龙贝格观测器（Luenberger Observer）或卡尔曼滤波器可以从含噪位置测量中估计速度，效果好于简单滤波。观测器利用电机动力学模型（通常是一阶模型：$\dot{\omega} = (K_t i - B\omega) / J$）与位置测量相结合，产生平滑的速度估计。我们会在第六章讲状态观测器时正式展开，这里先知道有这个选项即可——它是理论上的最优方案。

\subsubsection{编码器类型与选型}

并不是所有编码器都一样。\textbf{增量式编码器}（Incremental Encoder）输出两路正交脉冲（A、B 通道），可选一个每转一次的索引脉冲（Z 通道）。正交编码既给方向也给位置，大多数微控制器的定时器外设可以零 CPU 开销地硬件解码正交信号。增量式编码器断电后丢失位置，上电需要做归零操作。

\textbf{绝对式编码器}（Absolute Encoder）任何时候都能报告精确的轴角度，即使掉电重启也不丢。它们使用光学或磁感应原理，数字输出（SPI、SSI 或模拟量）。AS5048A 和 AS5600 是机器人圈子里常用的磁绝对式编码器，分辨率一般 12--14 位（$4096$--$16384$ 计数/转），非常优秀。缺点是通信延迟（SPI 几微秒）在极高速控制回路中可能有影响，而且有些廉价磁编码器存在 $\pm 1$--$2$ LSB 的非线性，随轴角变化。

对大多数比赛机器人来说，在电机轴（减速箱前）装一个 1024--4096 线的增量式编码器，是分辨率、速度和简洁性的最佳组合。如果需要绝对位置（比如关节上电后无需归零就要知道角度），就在输出轴装一个磁绝对式编码器。

%-----------------------------------------------------------------------------
\subsection{磁力计标定}
%-----------------------------------------------------------------------------

磁力计测量地球磁场来确定航向。在理想世界里，把传感器在水平面内转 $360\textdegree{}$，$x$-$y$ 读数会画出一个以原点为圆心的完美圆。现实中你得到的是一个偏了心的椭圆——标定就是来修这个的。

\textbf{硬铁偏差}（Hard Iron）来自传感器附近的永磁体或被磁化的铁磁材料（钢螺丝、扬声器磁铁、电池触点）。它们产生恒定磁场，让圆心偏离原点。修正方法：找到偏移圆的圆心，减掉——那就是你的硬铁偏移向量。

\textbf{软铁偏差}（Soft Iron）来自磁导率较高的材料（软钢、铁），它们非均匀地扭曲地磁场，把圆拉成椭圆。修正方法：对数据拟合椭圆，计算出把椭圆变回圆的变换矩阵。

实际标定流程：让机器人慢慢转 $360\textdegree{}$（更好的做法是各方向翻滚），记录磁力计数据，对 $x$-$y$ 散点图拟合椭圆，提取偏移和缩放参数。很多开源工具可以自动完成。

一个忠告：磁力计在室内\textbf{不可靠}。钢结构梁、钢筋混凝土、电线和附近的电机，都会以不可预测、位置相关的方式扭曲磁场。如果比赛场地在钢框架建筑里，你的磁力计航向可能在场地不同位置偏差 $10$--$30\textdegree{}$。用磁力计做粗略的航向初始化可以，但别指望它做室内连续导航。对于几分钟的短比赛，陀螺仪积分（加上周期性修正）几乎总是更靠谱。

%-----------------------------------------------------------------------------
\subsection{电流传感器标定}
%-----------------------------------------------------------------------------

如果你在做无刷电机的磁场定向控制（FOC, Field-Oriented Control），电流测量直接进入最内层控制回路，这里的误差会传遍全局。

\textbf{零偏标定：}电机启动前、电流为零时，读几百个电流传感器采样取平均，这就是零电流偏移，后续所有读数都要减掉它。很多电机驱动板上，电流传感器是分流电阻加运放，参考电压在 $V_{cc}/2$ 附近，所以"零"读数大约是 $V_{cc}/2$，而不是真正的 $0\;\mathrm{V}$。

\textbf{带宽很重要。}电流环一般跑 $10$--$20\;\mathrm{kHz}$。如果电流传感器的带宽只有 $5\;\mathrm{kHz}$，传感器本身就变成了一个低通滤波器，在控制器需要作用的频率上引入显著相位滞后。这个相位滞后会吃掉相位裕度，可能导致电流环失稳。查传感器数据手册：带宽至少应该是电流环频率的 $5\times$。霍尔效应电流传感器（如 ACS712）带宽一般几百 kHz，够用。分流电阻加快速运放可以做到 MHz 级带宽，绰绰有余。

%-----------------------------------------------------------------------------
\subsection{传感器融合架构}
%-----------------------------------------------------------------------------

传感器标定好了，下一个问题是怎么把它们组合起来。有两种根本不同的思路。

\textbf{测量级融合}（Measurement-level Fusion）直接混合原始（标定后）的传感器数据。第二章的互补滤波器就是经典案例：低通加速度计、高通陀螺仪、加到一起。简单轻量，在你恰好有两个噪声特性互补的传感器时非常好使。

\textbf{状态级融合}（State-level Fusion）利用系统动力学模型，从所有传感器数据中同时估计系统的内部状态（位置、速度、姿态）。卡尔曼滤波器是这方面的黄金标准。它能处理任意数量的传感器，分别考虑各传感器的噪声特性，并在一定假设下给出统计最优的估计。

怎么选？一条简单的判断规则：

如果你有 1--2 个传感器、噪声特性互补且关系清楚，就用互补滤波器。五行代码，微秒级运算，调参容易。

如果你有 3 个或更多传感器，或者测量与状态之间的关系是非线性的（比如 GPS + IMU + 磁力计做全姿态估计），或者你想要可证明的最优性能，就用第六章的扩展卡尔曼滤波器（EKF）。EKF 在统一框架下处理所有这些情况。代价是更多代码、更多调参、更多运算量——但现代微控制器轻松应付。

大多数比赛机器人属于互补滤波器阵营：加速度计-陀螺仪互补滤波做倾角估计，可能再加上编码器里程计做位置。EKF 留给真正需要的时候。

一个好用的心智模型：把传感器融合想成\textit{频域分工}。加速度计在低频段可信（看到的是恒定不变的重力），但高频被振动污染。陀螺仪在高频段可信（对旋转瞬间响应），但低频会漂（偏置累积）。互补滤波器就是把每个传感器分配到它擅长的频段。卡尔曼滤波器做的是同一件事，只不过它根据噪声统计量自动计算最优交叉频率。

\subsection{传感器选型指南}

标定之前，你得先选对传感器。下面是一份针对比赛机器人常见感知需求的简明选型指南。

\textbf{姿态（倾斜/航向）：}六轴 IMU（加速度计 + 陀螺仪）覆盖大部分需求。只有在需要绝对航向、且场地附近没有大型钢结构时才加磁力计。BMI088、ICM-42688、MPU-6500 是常见选择，按品质递增排列。

\textbf{轮速：}使用至少 $N = 4096$ 线的正交编码器，才能在 $1\;\mathrm{kHz}$ 采样下平滑测到 $\sim\!5\;\mathrm{RPM}$ 的速度。如果电机有减速箱，在减速箱前的电机轴上计数，分辨率更高。

\textbf{距离/测距：}飞行时间（ToF, Time-of-Flight）传感器如 VL53L1X，毫米级分辨率，量程 $\sim\!4\;\mathrm{m}$。体积小、便宜、I2C 接口用起来简单。更远距离可以用 2D 激光雷达（如 RPLIDAR A1），$360\textdegree{}$ 扫描、$5$--$10\;\mathrm{Hz}$、量程 $12\;\mathrm{m}$。

\textbf{FOC 电流检测：}用分流电阻加快速运放，带宽最优。霍尔效应传感器（ACS712、ACS714）更容易集成，但带宽低、噪声大。电流传感器带宽至少应是电流环频率的 $5\times$。

总原则：买你负担得起的最好的传感器，然后认真标定。一颗标定良好的 3 块钱传感器，能吊打一颗标定稀烂的 30 块钱传感器。

%-----------------------------------------------------------------------------
\subsection{实践：验证标定质量}
%-----------------------------------------------------------------------------

标定不是"设完拉倒"的事。你需要验证它确实生效了，还需要记录正确的数据以便事后排查问题。

\textbf{加速度计健全性检查：}标定后，把 IMU 摆到几个任意朝向，计算 $\|\mathbf{a}_{\text{cal}}\|$，在每个朝向都应该是 $9.81 \pm 0.05\;\mathrm{m/s^2}$。如果偏了，说明标定矩阵有问题——多半是六面标定时有某个朝向没放平。

\textbf{陀螺仪健全性检查：}机器人静止时，对标定后的陀螺仪速率积分 60 秒，累积角度应在 $\pm 1\textdegree{}$ 以内。如果漂移更大，说明偏置估计不准。重新做静态标定，确保机器人真的一动不动（没有风扇引起振动、没有人碰桌子）。

\textbf{磁力计健全性检查：}做完硬铁/软铁标定后，在水平面内转 $360\textdegree{}$，画标定后 $x$-$y$ 数据的散点图。应该是一个以原点附近为圆心的圆，而不是椭圆或偏心块。

\textbf{该记录什么：}永远同时记录原始数据和标定后数据，带上时间戳。场上出了问题，你需要原始数据来判断是传感器本身出了毛病还是标定出了问题。如果只记了标定后的数据，你分不清。日志很便宜——每秒几 KB。省下的调试时间是巨大的。

\textbf{编码器健全性检查：}用手精确转轮子一整圈（在轮子上做个标记）。编码器应该精确报告 $N$ 个计数（或正交模式下 $4N$，取决于你的计数配置）。如果总是多一个或少一个，可能是电气噪声引起了假边沿——在编码器输出线上并一个小电容，或者缩短线缆长度。

\textbf{电流传感器健全性检查：}电机断开时，电流读数应该为零（减掉偏置后）。然后指令一个已知的恒定电流，用万用表或钳形表验证。如果传感器读 $1.0\;\mathrm{A}$ 但万用表说 $0.95\;\mathrm{A}$，说明你的增益标定偏了 5\%。

\textbf{何时重新标定：}陀螺仪偏置每次上电都要重新标——它在不同次开机之间会漂移几十分之一度每秒。加速度计和磁力计的标定只要不改安装位置、不新增铁磁部件就是稳定的。电流传感器偏置应该在每次启动时、电机通电之前测量。

%-----------------------------------------------------------------------------
% 案例分析
%-----------------------------------------------------------------------------

\vspace{0.5em}
\begin{mdframed}[linewidth=1pt, innertopmargin=10pt, innerbottommargin=10pt]
\textbf{案例分析：航向 5 分钟漂了 $30\textdegree{}$}

\vspace{0.3em}
问题：一台比赛机器人用磁力计做航向估计。测试中，航向在 5 分钟内缓慢漂移了 $30\textdegree{}$，尽管机器人已经回到了起始朝向。

排查从画原始磁力计 $x$-$y$ 数据的散点图开始——让机器人转 $360\textdegree{}$。结果不是以原点为圆心的圆，而是圆心大约在 $(150, -80)\;\mu\mathrm{T}$ 处。这是典型的硬铁偏差——钢制底盘和驱动电机里的钕磁铁产生了恒定磁偏移。

修复方法：让机器人慢慢转一整圈 $360\textdegree{}$（在平面上），记录磁力计数据，算出圆心 $(c_x, c_y)$，然后从所有后续读数中减掉：

\begin{verbatim}
mag_cal_x = mag_raw_x - cx;
mag_cal_y = mag_raw_y - cy;
heading = atan2(mag_cal_y, mag_cal_x);
\end{verbatim}

三行代码修完后，5 分钟航向误差从 $30\textdegree{}$ 降到不到 $3\textdegree{}$。剩余误差来自软铁偏差和室内磁场的非均匀性——对于还有陀螺仪积分兜底的比赛机器人来说，完全可以接受。

教训：在信任磁力计之前，永远先把数据画成散点图看一眼。圆心不在原点？硬铁问题。不是圆而是椭圆？还有软铁问题。
\end{mdframed}

\vspace{0.5em}
\begin{mdframed}[linewidth=1pt, innertopmargin=10pt, innerbottommargin=10pt]
\textbf{案例分析：速度在 10 RPM 以下疯了}

\vspace{0.3em}
问题：一台轮式机器人使用 1024 线正交编码器做轮速反馈。PID 速度控制器在 $60\;\mathrm{RPM}$ 以上表现良好，但低于约 $10\;\mathrm{RPM}$ 时电机会发出明显的顿挫声，速度曲线看起来像方波，在 $0$ 和某个最小值之间跳来跳去。

分析：$1\;\mathrm{kHz}$ 采样下，1024 线编码器的最小可检测速度为：

\[
v_{\min} = \frac{2\pi}{1024 \times 0.001\;\mathrm{s}} \approx 6.14\;\mathrm{rad/s} \approx 58.6\;\mathrm{RPM}
\]

低于这个速度，位置差 $\theta[n] - \theta[n-1]$ 几乎总是 0（没数到边沿），偶尔跳到 $\pm 1$ 个计数。差分速度在 0 和 $v_{\min}$ 之间振荡。PID 控制器看到"零速度"就猛加输出；然后看到"$v_{\min}$"又赶紧减。这就造成了听得见的顿挫。

修复结合了两个改动。第一，把速度计算频率降到每 $10\;\mathrm{ms}$ 一次（而非每 $1\;\mathrm{ms}$），$v_{\min}$ 降到 $5.86\;\mathrm{RPM}$。第二，加一个一阶低通滤波（$\alpha = 0.1$）来平滑剩余的量化台阶：

\begin{verbatim}
// Every 10 ms:
v_raw = (encoder_count - prev_count) * (2*PI/1024) / 0.01;
v_filt += 0.1 * (v_raw - v_filt);
prev_count = encoder_count;
\end{verbatim}

改完后电机平稳运行到大约 $3\;\mathrm{RPM}$。更低速度下，团队后来改用了 STM32 定时器输入捕获模式的 M/T 方法，直接测量编码器边沿间的时间，从而在接近零速时也能提供精细的速度分辨率。

教训：在造机器人\textit{之前}，先算好你的编码器-采样率组合对应的 $v_{\min}$。如果 $v_{\min}$ 高于你需要的最低速度，要么用更高分辨率的编码器，要么降低速度采样率，要么上 M/T 测量。
\end{mdframed}

\vspace{0.5em}
\subsection{标定清单}

下面是我们建议在每次组装新机器人或做了重大硬件变更后执行的标定流程。大约 30 分钟，能省掉后面几个小时的排查。

\begin{center}
\begin{tabular}{rll}
\toprule
\textbf{步骤} & \textbf{传感器} & \textbf{流程} \\
\midrule
1 & 电流传感器 & 零电流偏置（电机断开） \\
2 & 陀螺仪 & 60 秒静态偏置（机器人静置于平面） \\
3 & 加速度计 & 六面静态标定（首次或重新安装时） \\
4 & 磁力计 & $360\textdegree{}$ 旋转做硬铁/软铁标定（如果使用） \\
5 & 编码器 & 转一圈计数检查（验证精确 $N$ 个计数） \\
6 & 所有传感器 & 记录原始 + 标定数据，运行健全性检查 \\
\bottomrule
\end{tabular}
\end{center}

步骤 1 和 2 每次上电都要做。步骤 3--5 只有硬件变更时才需重做。步骤 6 应该养成习惯——花五分钟验证标定，能让你在比赛时免于追幽灵 Bug。

这整章的核心就一句话：\textit{先标定再滤波，先滤波再控制。}一个说实话的传感器，即使带着噪声，也比一个安安静静说谎的传感器有用无数倍。先拿到真相；本书后面会教你怎么用好它。
